\documentclass[../main.tex]{subfiles}
\begin{document}

% Nguồn: PLAN.md mục 5.1 và src/thesis_sumo/schema_observed.py (lược đồ đã cài đặt, quyết định R3, R9, R11).
\section{Lược đồ dữ liệu quan sát}
\label{section:appendix_schema}

Lược đồ dữ liệu quan sát quy định các tệp mà mỗi ngày dữ liệu phải có. Bộ sinh dữ liệu thay thế PROXY\_D1--PROXY\_D5 và bộ chuyển đầu ra mô phỏng cùng ghi theo lược đồ này, và dữ liệu trích từ video phải tuân theo nó; nhờ đó một bộ hàm chỉ số duy nhất áp dụng được cho cả dữ liệu quan sát lẫn dữ liệu mô phỏng (Mục~\ref{section:3.6}). Danh sách cột được định nghĩa tại một mô-đun duy nhất trong mã nguồn để các thành phần không giữ các bản sao lệch nhau. Bảng~\ref{table:appendix_schema} trình bày lược đồ đã cài đặt.

\begin{table}[!htbp]
\centering
\small\setlength{\tabcolsep}{3pt}
\caption{Lược đồ dữ liệu quan sát đã cài đặt}
\label{table:appendix_schema}
\begin{tabular}{|p{3.7cm}|p{5.1cm}|p{4.2cm}|}
\hline
\textbf{Tệp} & \textbf{Cột} & \textbf{Nội dung} \\ \hline
\texttt{turn\_\allowbreak counts.csv} & \texttt{interval\_\allowbreak start}, \texttt{interval\_\allowbreak end}, \texttt{approach}, \texttt{movement}, \texttt{vclass}, \texttt{count} & Số đếm theo ô 15 phút (900~s); \texttt{movement} thuộc \{\texttt{l}, \texttt{s}, \texttt{r}, \texttt{u}, \texttt{illegal\_\allowbreak l}\} \\ \hline
\texttt{signal\_\allowbreak events.csv} & \texttt{t}, \texttt{approach\_\allowbreak group}, \texttt{state} & Thời điểm đổi trạng thái đèn; \texttt{state} thuộc \{\texttt{G}, \texttt{y}, \texttt{r}\} \\ \hline
\texttt{stopline\_\allowbreak crossings.csv} & \texttt{t}, \texttt{approach}, \texttt{lane\_\allowbreak or\_\allowbreak lat\_\allowbreak m}, \texttt{vclass}, \texttt{track\_\allowbreak id}, \texttt{signal\_\allowbreak state\_\allowbreak at\_\allowbreak t} & Mỗi lượt xe qua vạch dừng, kèm trạng thái đèn tại thời điểm đó \\ \hline
\texttt{queue\_\allowbreak snapshots.csv} & \texttt{cycle\_\allowbreak id}, \texttt{t\_\allowbreak green\_\allowbreak onset}, \texttt{approach}, \texttt{queue\_\allowbreak length\_\allowbreak m}, \texttt{n\_\allowbreak car}, \texttt{n\_\allowbreak moto}, \texttt{area\_\allowbreak m2} & Hàng chờ tại thời điểm bắt đầu đèn xanh của mỗi chu kỳ \\ \hline
\texttt{lateral\_\allowbreak snapshots.csv} & \texttt{cycle\_\allowbreak id}, \texttt{approach}, \texttt{track\_\allowbreak id}, \texttt{vclass}, \texttt{lat\_\allowbreak m}, \texttt{long\_\allowbreak m\_\allowbreak from\_\allowbreak stopline} & Vị trí ngang và dọc của xe đang dừng chờ \\ \hline
\texttt{box\_\allowbreak queue.csv} & \texttt{cycle\_\allowbreak id}, \texttt{t}, \texttt{approach}, \texttt{n\_\allowbreak waiting\_\allowbreak in\_\allowbreak box}, \texttt{t\_\allowbreak clear\_\allowbreak start}, \texttt{t\_\allowbreak clear\_\allowbreak end} & Xe rẽ trái trái phép chờ giữa nút, mỗi dòng ứng với một lần kết thúc pha xanh đi thẳng của hướng đối diện \\ \hline
\texttt{traversals.csv} & \texttt{track\_\allowbreak id}, \texttt{vclass}, \texttt{approach}, \texttt{movement}, \texttt{exit\_\allowbreak approach}, \texttt{lat\_\allowbreak m\_\allowbreak at\_\allowbreak stopline}, \texttt{t\_\allowbreak range\_\allowbreak entry}, \texttt{t\_\allowbreak stopline}, \texttt{t\_\allowbreak exit} & Bảng dẫn xuất, mỗi dòng một lượt qua vạch dừng; phục vụ M6 (thời gian qua nút) và M8 (đi sai làn) \\ \hline
\texttt{queue\_\allowbreak timeseries.csv} & \texttt{t}, \texttt{approach}, \texttt{queue\_\allowbreak length\_\allowbreak m}, \texttt{n\_\allowbreak car}, \texttt{n\_\allowbreak moto} & Hàng chờ của mỗi nhánh theo lưới thời gian 2~s, kể cả khi không có hàng chờ; cùng định nghĩa với \texttt{queue\_\allowbreak snapshots.csv} \\ \hline
\texttt{space\_\allowbreak seeking\_\allowbreak events.csv} & \texttt{t}, \texttt{cycle\_\allowbreak id}, \texttt{approach}, \texttt{track\_\allowbreak id}, \texttt{vclass}, \texttt{space}, \texttt{queue\_\allowbreak length\_\allowbreak m\_\allowbreak at\_\allowbreak t} & Mỗi dòng một cặp (xe, không gian) ở lần đầu xe được thấy trong không gian đó; \texttt{space} thuộc \{\texttt{wrong\_\allowbreak lane}, \texttt{right\_\allowbreak lane}, \texttt{sidewalk}, \texttt{opposite}\} \\ \hline
\texttt{right\_\allowbreak turn\_\allowbreak delay.csv} & \texttt{track\_\allowbreak id}, \texttt{approach}, \texttt{vclass}, \texttt{t\_\allowbreak enter}, \texttt{t\_\allowbreak exit}, \texttt{blocked} & Mỗi dòng một lượt rẽ phải qua vạch dừng từ lòng đường; \texttt{blocked} cho biết có xe đi thẳng dừng trên làn 0 phía trước \\ \hline
\end{tabular}
\end{table}

\begin{sloppypar}
Các quy ước chung của lược đồ như sau. Mọi thời điểm tính bằng giây kể từ đầu cửa sổ quan sát (phía mô phỏng là giây mô phỏng). Mã nhánh \texttt{approach} là một trong \texttt{LVL\_W}, \texttt{LVL\_E}, \texttt{HMG}, \texttt{NT} (trên N1--N6, hai nhánh phụ mang mã \texttt{HDT\_N}, \texttt{HDT\_S}); \texttt{approach\_group} trong \texttt{signal\_events.csv} cũng là mã nhánh và biểu thị trạng thái đèn của hướng đi thẳng của nhánh đó; các hướng có mũi tên riêng (rẽ trái, quay đầu của Lê Văn Lương) không có nhóm trong tệp này. \texttt{signal\_state\_at\_t} trong \texttt{stopline\_crossings.csv} là trạng thái liên kết của chính xe qua vạch. Trạng thái đèn được chuẩn hóa về ba giá trị: \texttt{G} (xanh, gồm cả trạng thái xanh phải nhường của SUMO), \texttt{y} (vàng) và \texttt{r} (mọi trạng thái buộc dừng còn lại). \texttt{vclass} nhận \texttt{car} hoặc \texttt{moto}; xe máy vi phạm được tính là \texttt{moto}. Chu kỳ được xác định theo từng nhánh từ các thời điểm bắt đầu đèn xanh: chu kỳ thứ $k$ là khoảng giữa lần bắt đầu xanh thứ $k$ và thứ $k+1$, còn các thời điểm trước lần bắt đầu xanh đầu tiên thuộc chu kỳ $-1$ và bị loại khi tính chỉ số.
\end{sloppypar}

\begin{sloppypar}
Trong \texttt{box\_queue.csv}, \texttt{t} là thời điểm kết thúc pha xanh đi thẳng của hướng đối diện; \texttt{n\_waiting\_in\_box} là số xe rẽ trái trái phép đang dừng (tốc độ dưới 5~km/h) trong nút tại thời điểm đó; \texttt{t\_clear\_start} và \texttt{t\_clear\_end} là thời điểm rời nút đầu tiên và cuối cùng của tất cả xe vi phạm đã vào nút trong chu kỳ, kể cả xe rời sớm qua khoảng trống, nên \texttt{t\_clear\_start} có thể sớm hơn \texttt{t}. Trong \texttt{traversals.csv}, \texttt{exit\_approach} là nhánh xe đi ra, \texttt{t\_range\_entry} là thời điểm xe đi vào phạm vi quan sát hàng chờ phía trước vạch dừng, \texttt{t\_stopline} là thời điểm qua vạch và \texttt{t\_exit} là thời điểm xe vào làn ra. \texttt{lat\_m\_at\_stopline} là vị trí ngang của xe khi qua vạch dừng; M8 (đi sai làn) đọc làn của xe từ cột này (Mục~\ref{section:3.7}).
\end{sloppypar}

\begin{sloppypar}
Ba bảng cuối phục vụ nhóm tìm khoảng trống và chỉ suy từ vị trí xe, không đọc nhật ký can thiệp, nên đo được cả từ video. Trong \texttt{space\_seeking\_events.csv}, \texttt{right\_lane} là xe đi thẳng dừng (tốc độ dưới 5~km/h) trên làn 0, \texttt{sidewalk} là xe đi trên vỉa hè, và \texttt{wrong\_lane} là xe rẽ qua vạch dừng ở một làn không cho hướng của nó, theo đúng quy tắc của M8 (đi sai làn); giá trị \texttt{opposite} (đi ngược chiều) có trong lược đồ nhưng không được sinh ra, vì hành vi 8 nằm ngoài phạm vi mô hình. Trong \texttt{right\_turn\_delay.csv}, \texttt{t\_enter} là thời điểm xe vào phạm vi quan sát hàng chờ và \texttt{t\_exit} là thời điểm xe qua vạch dừng.
\end{sloppypar}

\begin{sloppypar}
Bảng \texttt{traversals.csv} là bảng dẫn xuất, do cả bộ chuyển đầu ra mô phỏng và bộ sinh dữ liệu thay thế cùng sinh ra. Lược đồ không có tầng quỹ đạo; các chỉ số chỉ dùng mười bảng trên.
\end{sloppypar}

\end{document}
